%% kapa: exam=1 round=20 year=2009 period=2 form=A subject=accounting subject_ko=회계학 seq=11 number=11
%% source: Q-Net 감정평가사 기출문제 (cst00309) artlSeq=5024178 file=extracted/1008/5024178/2009년 제20회 감정평가사 1차시험 기출문제/2교시.hwp
%% conversion: kordoc 4.24.51 -> tools/gen1.py (auto); review=proofread-sonnet; flags=none
\begin{question}{11}
다음 중 회계변경의 회계처리에 대한 내용으로 옳지 \uline{않은} 것은?
\begin{choices}
  \item 중요성의 판단에 따라 기업회계기준과 다르게 회계처리하던 항목들의 중요성이 커지게 되어 기업회계기준을 적용하는 경우는 회계정책의 변경으로 보고 회계처리한다.
  \item 단순히 세법의 규정을 따르기 위한 회계변경은 정당한 회계변경으로 보지 않는다.
  \item 정당한 회계정책변경의 경우 회계변경에 따른 누적효과를 합리적으로 결정하기 어렵다면 전진적으로 처리할 수 있다.
  \item 정당한 회계변경의 경우 회계변경의 속성상 그 효과를 회계정책의 변경효과와 회계추정의 변경효과로 구분하기 불가능하다면 이를 회계추정의 변경으로 본다.
  \item 정당한 회계정책 변경의 경우 비교목적으로 공시되는 전기 재무제표는 소급적용에 따른 수정사항을 반영하여 재작성한다.
\end{choices}
\end{question}
\answer{1}
